\documentclass[11pt, a4paper]{article}
% ---- Core Geometry & Margins ----
\usepackage[margin=1in, top=1.2in, bottom=1.2in]{geometry}
\setlength{\parindent}{0pt}
\setlength{\parskip}{0.8em}
% ---- Font & Micro-Typography ----

\usepackage[utf8]{inputenc}
\usepackage[T1]{fontenc}
\usepackage{lmodern}
% \usepackage{newtxtext}
% \usepackage{newtxmath}
% \usepackage{libertinus}
\usepackage{microtype}

% ---- Language & Math ----
\usepackage[english]{babel}
\usepackage{amsmath, amssymb, amsthm}
\usepackage{booktabs}  % In your preamble

% ---- Formatting & Aesthetics ----
\usepackage{booktabs}
\usepackage{graphicx}
\usepackage{xcolor}
\usepackage{csquotes}

% ---- Color Palette ----
\definecolor{navyblue}{RGB}{16, 55, 87}
\definecolor{darkslate}{RGB}{53, 64, 83}

% ---- Hyperlinks Setup ----
\usepackage[
  pdfencoding=auto,
  colorlinks=true,
  linkcolor=navyblue,
  citecolor=navyblue,
  urlcolor=navyblue
]{hyperref}

% ---- Bibliography Setup ----
% \usepackage[
%   backend=biber,
%   style=numeric-comp,
%   sorting=nyt,
%   giveninits=true
% ]{biblatex}
% \addbibresource{references.bib}

% --- Metadata ---
\title{Random Thoughts}
\author{Guetchuin Tadouanla}
\date{today}


% ====================================================================================
\begin{document}
\pagenumbering{roman}
\setcounter{page}{1}

\tableofcontents
\newpage
\pagenumbering{arabic}
\setcounter{page}{1}

\section{Introduction}
According to the \textit{Sustainable Development Report 2025} \cite{sdr2025}, Cameroon ranks $\mathbf{135^{th}}$ out of \textbf{169} countries, with an SDG Index score of $\mathbf{57.6\%}$. Figure~\ref{fig:SDG_dashboard} displays the overall national status across all 17 goals.

\begin{figure}[htbp]
  \centering
  \includegraphics[width=0.95\textwidth]{figures/SDG_dashboard.png}
  \caption{Cameroon SDG Dashboard and Trends}
  \label{fig:SDG_dashboard}
\end{figure}

Cameroon has achieved \textbf{SDG 12 (Responsible Consumption and Production)} and \textbf{SDG 13 (Climate Action)}. \textbf{SDG 14} retains moderate challenges, while \textbf{SDG 8} faces significant challenges. However, the remaining SDGs face \textbf{major challenges} with stagnating progress. This report focuses on \textbf{SDG 7 (Affordable and Clean Energy)}, \textbf{SDG 9 (Industry, Innovation, and Infrastructure)}, and \textbf{SDG 11 (Sustainable Cities and Communities)}, outlining targeted embedded and IoT engineering solutions.

\section{SDG 7: Affordable and Clean Energy}
SDG 7 aims to ensure access to affordable, reliable, sustainable, and modern energy for all. Cameroon faces major challenges, with progress stagnating or advancing at under $50\%$ of the required pace. Key tracking indicators include:

\begin{itemize}
  \itemsep 0em
  \item Population with access to electricity
  \item Population with access to clean fuels and technology for cooking
  \item $\text{CO}_2$ emissions from fuel combustion per total electricity output
  \item Renewable energy share in total final energy consumption
\end{itemize}

\begin{figure}[htbp]
  \centering
  \includegraphics[width=0.85\textwidth]{figures/sustainable-development-report-2026-population-with-access-to-clean-fuels-and-technology-for-cooking.png}
  \caption{Population with access to clean fuels and technology for cooking}
  \label{fig:sdg7_cooking}
\end{figure}

\subsection{Indicator: Population with Access to Clean Fuels and Technology for Cooking}

\textbf{Problem Statement:} Over $80\%$ of low-income Cameroonian households rely on firewood and charcoal, causing severe indoor air pollution and health risks. High upfront costs ($15\text{--}20~\text{USD}$) for Liquefied Petroleum Gas (LPG) cylinders prevent adoption compared to purchasing cheap daily charcoal.

\textbf{Potential Solutions:}
\begin{enumerate}
  \itemsep 0em
  \item \textbf{Smart PAYGO LPG Systems \cite{shupler2021}:} Low-cost MCU smart meters with load-cell weight sensors, electronic solenoid valves, and cellular modems (NB-IoT/GSM). Users pay small amounts via mobile money (MTN MoMo/Orange Money), unlocking gas incrementally while notifying suppliers before empty.
  \item \textbf{IoT Biogas Digesters:} Low-power sensor nodes ($pH$, pressure, temp) powered by cookstove waste heat via Thermoelectric Generators (TEGs) to monitor organic waste digestion in livestock regions, reporting telemetry over LoRaWAN.
  \item \textbf{Thermoelectric Gasifiers:} Forced-draft micro-gasifiers using Seebeck TEG modules to power PWM combustion fans directly from stove heat, ensuring clean biomass burning off-grid without electricity.
\end{enumerate}

\section{SDG 9: Industry, Innovation, and Infrastructure}
SDG 9 aims to build resilient infrastructure, promote inclusive industrialization, and foster innovation. Cameroon faces major challenges, with progress stagnating or advancing at under $50\%$ of the required pace. Key tracking indicators include:

\begin{itemize}
  \itemsep 0em
  \item Rural population with access to all-season roads
  \item Mobile broadband subscriptions per 100 population
  \item Logistics Performance Index: Infrastructure Score
  \item Expenditure on research and development (\% of GDP)
  \item Total patent applications by applicant's origin
\end{itemize}

\subsection{Indicator: Total Patent Applications by Applicant's Origin}

\textbf{Problem Statement:} Domestic patent applications remain extremely low, reflecting a dependency on expensive, imported technology unsuited for local tropical environments. Local engineers lack hardware R\&D funding, prototyping labs, and commercialization pathways.

\textbf{Potential Solutions:}
\begin{enumerate}
  \itemsep 0em
  \item \textbf{Frugal Embedded R\&D:} Local engineering focused on open microcontrollers (STM32/ESP32) and modular PCBs tailored for agricultural processing, energy metering, and environmental monitoring.
  \item \textbf{Open-Source Hardware Ecosystems:} Utilizing open architectures (FreeRTOS, Linux, open gateway designs) to allow local innovators to build and deploy systems without prohibitive licensing fees.
  \item \textbf{University FabLabs:} Establishing rapid prototyping labs (3D printers, PCB mappers) in universities to transition student designs into bench-tested, patentable hardware.
\end{enumerate}

\section{SDG 11: Sustainable Cities and Communities}
SDG 11 aims to make human settlements inclusive, safe, resilient, and sustainable. Cameroon faces major challenges, with progress stagnating or advancing at under $50\%$ of the required pace. Key tracking indicators include:

\begin{itemize}
  \itemsep 0em
  \item Proportion of urban population living in slums
  \item Access to an improved water source, piped (\% of urban population)
  \item Satisfaction with public transport
  \item Annual mean concentration of particulate matter ($\text{PM}_{2.5}$) in urban areas
\end{itemize}

\subsection{Indicator: Access to an Improved Water Source, Piped}

\textbf{Problem Statement:} Municipal distribution networks in Douala and Yaoundé suffer severe pipe leaks and supply outages, leaving urban populations without clean piped water and forcing reliance on untreated sources or expensive vendors.

\textbf{Potential Solutions:}
\begin{enumerate}
  \itemsep 0em
  \item \textbf{Smart Pressure/Acoustic WSNs:} Low-power LoRaWAN/NB-IoT sensor nodes monitoring acoustic vibrations and pressure drops along water mains to detect pipe bursts in real time.
  \item \textbf{Prepaid IoT Water Kiosks:} Solar-powered neighborhood kiosks with RFID readers, MCU flow meters, and mobile money billing to dispense clean piped water at fair municipal rates.
  \item \textbf{Autonomous Water Quality Telemetry:} Solar-powered purification stations with real-time $pH$, turbidity, and conductivity sensors broadcasting drinking safety alerts.
\end{enumerate}

\section{Conclusion}
Progress across SDGs 7, 9, and 11 in Cameroon can be accelerated through frugal embedded engineering. Deploying smart PAYGO LPG meters, open-source hardware, and low-power water monitoring networks offers low-cost, locally maintainable solutions to national sustainability bottlenecks.

\newpage
\begin{thebibliography}{9}

  \bibitem{sdr2025}
  J.~Sachs, G.~Lafortune, and G.~Fuller, \textit{Sustainable Development Report 2025}. Sustainable Development Solutions Network (SDSN), 2025.

  \bibitem{shupler2021}
  M.~Shupler et al., ``Pay-as-you-go LPG for clean cooking in low-income urban settings: a community-based trial,'' \textit{Nature Energy}, vol. 6, no. 4, pp. 379--389, 2021.

\end{thebibliography}

\end{document}
